\documentclass{article}

\usepackage[a4paper]{geometry}
\usepackage[utf8]{inputenc}
\usepackage[english]{babel}
\usepackage[noend]{algpseudocode}
\usepackage{amsthm}
\usepackage{colortbl}
\usepackage{graphicx}
\usepackage{lineno}
\usepackage[textwidth=4cm]{todonotes}

\geometry{left=3cm, right=6cm, top=3cm, bottom=3cm, marginparwidth=4cm} 

\linenumbers

\setuptodonotes{noline, color=green!15, size=\scriptsize}

\newtheorem*{exercise}{Exercise}
\newcommand{\grey}{\cellcolor{black!25}}

\title{Algorithms and Data Structures \\
  Handin 0: Pair sum}

\author{
    202605673 Lukas Alstrup \\
    202605116 Thomas Seeberg Hansen \\
    202606296 Frederik Dommer Løvbjerg Plum
}

%\date{September 1, 2025}  % Use \date to set a specific date

\begin{document}

\maketitle
\begin{exercise}[Squared]
	Prove that the algorithm Squared(n) computes $s = n^2$. Find an appropriate loop invariant I, and use this to conclude that $s = n^2$ when the algorithm terminates.
\end{exercise}

\section*{Comment}
\section*{Simple solution}
$$s_x = 2\sum_{i=0}^{x}i - x$$
$$s_x = 2 \cdot \frac{x(x+1)}{2}-x$$
$$s_x = x \cdot (x+1) -x$$
$$ s_x = x^2 + x - x$$
$$ s_x = x^2$$

\begin{equation}
	0 \leq i \leq n \land 0 \leq s \leq n^2
\end{equation}

\section*{Efficient solution}
\end{document}
